% Part: lambda-calculus
% Chapter: introduction
% Section: reduction

\documentclass[../../../include/open-logic-section]{subfiles}

\begin{document}

\olfileid[fr]{lam}{int}{red}
\olsection{Réduction des termes lambda}

Que peut-on faire avec les termes lambda?
Les simplifier. Si $M$ et $N$ sont des termes
lambda quelconques et $x$ une variable,
$\Subst{M}{N}{x}$ désigne le résultat de la
substitution de~$N$ aux occurrences libres de~$x$ dans~$M$, après
avoir renommé les variables liées de~$M$
qui risqueraient de capturer des variables
libres de~$N$. Par exemple,
\[
\Subst{(\lambd[w][xxw])}{yyz}{x}
=\lambd[w][(yyz)(yyz)w].
\]

\begin{digress}
On rencontre aussi les notations $[N/x]M$,
$[x/N]M$ et $M[x/N]$ pour la substitution.
Attention aux conventions!
\end{digress}

Intuitivement, $(\lambd[x][M])N$ et
$\Subst{M}{N}{x}$ ont le même sens. Le
remplacement du premier terme par le second
s'appelle une \emph{contraction~$\beta$}.
Le terme $(\lambd[x][M])N$ est un
\emph{rédex} et $\Subst{M}{N}{x}$ son
\emph{contractum}. Plus généralement, si
l'on peut transformer un terme~$P$ en~$P'$
par contraction~$\beta$ d'un de ses
sous-termes, nous disons que $P$ se
\emph{réduit en une étape} à~$P'$, et nous
écrivons $P\redone P'$. Si l'on peut
obtenir~$P'$ à partir de~$P$ par un nombre
quelconque de telles étapes, éventuellement
nul, alors $P$ se \emph{réduit} à~$P'$;
en symboles, $P\red P'$. Un terme que l'on
ne peut plus réduire par contraction~$\beta$
est dit \emph{$\beta$-irréductible}, ou
\emph{$\beta$-normal}. Lorsque le contexte
est clair, nous dirons simplement « se
réduit » plutôt que « se réduit par~$\beta$ »,
et de même pour les termes apparentés.

Considérons quelques exemples.
\begin{enumerate}
\item Nous avons
\begin{align*}
(\lambd[x][xxy]) \lambd[z][z] & \redone (\lambd[z][z])(\lambd[z][z]) y \\
& \redone (\lambd[z][z]) y \\
& \redone y.
\end{align*}
\item « Simplifier » un terme peut le rendre
      plus complexe :
\begin{align*}
(\lambd[x][xxy])(\lambd[x][xxy]) & \redone (\lambd[x][xxy])(\lambd[x][xxy])y \\
& \redone (\lambd[x][xxy])(\lambd[x][xxy])yy \\
& \redone \dots
\end{align*}
\item Une réduction peut aussi laisser le
      terme inchangé :
\[
(\lambd[x][xx])(\lambd[x][xx])
\redone (\lambd[x][xx])(\lambd[x][xx]).
\]
\item Certains termes peuvent se réduire
      de plusieurs façons. Ainsi,
\[
(\lambd[x][(\lambd[y][yx]) z)] v
\redone (\lambd[y][yv]) z
\]
      en contractant l'application extérieure,
      tandis que
\[
(\lambd[x][(\lambd[y][yx]) z)] v
\redone (\lambd[x][zx]) v
\]
      résulte de la contraction de
      l'application intérieure. Dans ce cas,
      les deux termes obtenus se réduisent
      ensuite au même terme, $zv$.
\end{enumerate}

Le résultat final du dernier exemple n'est
pas dû au hasard : il illustre une propriété
profonde et importante du lambda-calcul,
appelée \emph{propriété de Church--Rosser}.

\end{document}
